% =============================================================
%  chapter4_abstract.tex
%  Shared abstract text for Chapter 4 (nGibbs). Kept separate from
%  the surrounding \begin{quote}...\end{quote} wrapper in
%  Chapter4/chapter4.tex for the same reason as the other chapters'
%  abstract files: the wrapper is class-specific, this raw text is not.
%  Numbers are written out rather than taken from the results macros,
%  because the abstract is typeset before those macros are \input.
% =============================================================
Thermodynamic models that compute equilibrium phase assemblages by free energy minimization, such as \textsc{MELTS} and
\textsc{HeFESTo}, are often the computational bottleneck of the workflows that use them. The problems that would benefit most from
these models include compositional inversions of seismic data, geodynamic models of a chemically heterogeneous mantle, and
physical models of realistic magmatic systems that are heterogeneous in composition and temperature over space and time. These problems 
require large numbers of equilibria over a composition space that lookup tables cannot cover. Recently, machine-learning surrogate models have 
been built to address this need, but those that have been built to date each emulate a narrowly defined slice of their parent model. This chapter presents
\emph{nGibbs}, a set of general deep-learning emulators that reproduce \textsc{MELTS} 1.0.2 and 1.2.0 and the \textsc{HeFESTo} models of
Earth's and Mars' mantles across broad ranges of pressure, temperature and bulk composition. Each emulator predicts the full
equilibrium assemblage, the abundance of every phase and the endmember composition of every solution phase. Like the parent models
\emph{nGibbs} emulates, phase abundances vary continuously across phase boundaries, and mass is conserved. Physical properties are computed from the
predicted assemblage with the parent models' own equations of state. On a validation dataset generated from bulk compositions distinct from the training data,
the emulators predict the presence of the major silicate and oxide phases with precision and recall of 0.93 or better and conserve every element to a median relative error of about
$5\times10^{-6}\%$. Along a depleted-mantle isentrope, emulated densities and seismic velocities differ from \textsc{HeFESTo} by
0.02\% on average. The \textsc{HeFESTo} emulators run four orders of magnitude faster than \textsc{HeFESTo} on an ordinary
CPU, and the \textsc{MELTS} emulators two orders of magnitude faster than \textsc{alphaMELTS}. The emulators are
differentiable, run on CPUs or GPUs, and are interchangable with \textsc{alphaMELTS} in existing \textsc{PetThermoTools} workflows. These fast, stable, and compositionally
general models make composition a practical free parameter in seismic inversions and coupled geodynamic models.
